% paper.tex — Flagging soiled rooftop PV from public data alone.
% Build: see README.md. Figures must be generated before compiling.
\documentclass[10pt,twocolumn]{article}

\usepackage[letterpaper,margin=0.85in]{geometry}
\usepackage[T1]{fontenc}
\usepackage[utf8]{inputenc}
\usepackage{times}
\usepackage{graphicx}
\usepackage{booktabs}
\usepackage{amsmath}
\usepackage{amssymb}
\usepackage{siunitx}
\usepackage{url}
\usepackage[numbers,sort&compress]{natbib}
\usepackage{caption}
\usepackage{subcaption}
\usepackage{xcolor}
\usepackage{microtype}
\usepackage[colorlinks=true,allcolors=blue!55!black]{hyperref}

\sisetup{detect-all,group-separator={,}}
\captionsetup{font=small,labelfont=bf}
\graphicspath{{build/}}

\newcommand{\artifact}[1]{{\small\texttt{#1}}}

% A heading that wraps is justified by default, which stretches its interword
% space across the whole column. Run-in \paragraph titles are the worst affected
% because they are the longest. Ragged right on all three levels.
\makeatletter
\renewcommand\section{\@startsection{section}{1}{\z@}%
  {-3.5ex \@plus -1ex \@minus -.2ex}{2.3ex \@plus.2ex}%
  {\normalfont\Large\bfseries\raggedright}}
\renewcommand\subsection{\@startsection{subsection}{2}{\z@}%
  {-3.25ex\@plus -1ex \@minus -.2ex}{1.5ex \@plus .2ex}%
  {\normalfont\large\bfseries\raggedright}}
% The negative afterskip in article's default makes \paragraph a RUN-IN heading:
% the title joins the body paragraph and is justified with it, so a title that
% nearly fills a line gets its interword space stretched across the column and
% \raggedright in the style above cannot reach it. These titles are sentences,
% not labels, so give them their own line (positive afterskip) instead.
\renewcommand\paragraph{\@startsection{paragraph}{4}{\z@}%
  {2.4ex \@plus1ex \@minus.2ex}{0.5ex \@plus .1ex}%
  {\normalfont\normalsize\bfseries\raggedright}}
\makeatother

\title{\bf Station Labels Cannot Rank Roofs:\\
Why Public-Data Soiling Models Do Not Transfer,\\
and What Cleaning Is Actually Worth}

\author{
  Cameron Gordon \quad Joshua Ramirez \quad Akshitha Nagaraj \quad Craig Fellers\\[2pt]
  {\normalsize Better Behavior Foundation}\\[2pt]
  {\small\texttt{betterbehaviorfoundation@gmail.com} \quad
   \url{https://betterbehaviorfoundation.com}}
}
\date{}

% Added when the abstract sat inline in column one, where twocolumn's default
% \flushbottom stretched the glue under it into a two-inch gap. The standalone title page
% removed that case, and the document is now byte-identical either way (13 pages, same 49
% badness warnings). Kept as a guard: the gap returns the moment any page runs short.
\raggedbottom

% Default float parameters are too tight for seven full-width floats; without this
% relaxation figures drift several pages past the text that references them.
\setcounter{topnumber}{2}
\setcounter{bottomnumber}{1}
\setcounter{totalnumber}{3}
\renewcommand{\topfraction}{0.92}
\renewcommand{\bottomfraction}{0.55}
\renewcommand{\textfraction}{0.08}
\renewcommand{\floatpagefraction}{0.75}
\renewcommand{\dbltopfraction}{0.92}
\renewcommand{\dblfloatpagefraction}{0.70}
\setcounter{dbltopnumber}{2}

\begin{document}

% Title, authors, and abstract stand alone on page 1 at full width; the body starts
% two-column on page 2. Switching column count only at a \clearpage keeps the float
% and footnote machinery from straddling the boundary.
\onecolumn
\maketitle
\thispagestyle{empty}

\begin{center}
\begin{minipage}{0.84\textwidth}
\begin{abstract}
\noindent
Whether a rooftop solar array is dirty enough to be worth cleaning is measurable only from
production telemetry, which vendors hold and homeowners, municipalities and service
providers cannot get. We ask how far a soiling-targeting system gets on public data alone,
and report where it fails and why.

Locating arrays is solved: tile-level box-$F_1$ \textbf{0.826} (95\% CI [0.798, 0.853]),
\SI{73.6}{\percent} recall against permits the detector never saw. Ranking them is not, and
our own validation hid that. Folds that hold out either \emph{site} or \emph{year} but not
both flatter the model: ours leaves the held-out year's sites in training for
\textbf{88.7}\% of rows. Closing the leak moves pooled AUC from 0.710 to \textbf{0.622}
(95\% CI [0.571, 0.670]); a size-matched control puts \SI{83}{\percent} of that fall on
leakage (95\% CI [52, 99]). Latitude and longitude alone score 0.644, and an unfitted 2006
empirical model scores 0.610.

The limit is the label unit. Public labels are measured at \emph{stations}, so every roof
in a catchment inherits one. On per-system telemetry, per-array features rank systems
sharing weather at \textbf{0.573} against the \textbf{0.5} a station-label model cannot
exceed by construction---but three standard cleaning assumptions put those systems' median
annual loss at 2.72, 9.80 and \SI{20.93}{\percent} and rank them at $\rho$ 0.24 to 0.89.
Per-system labels are necessary and not yet sufficient.

A third result reorders the first two, and is immune to the second. On 149 metered
rooftops one wash recovers a median \textbf{\$28.10} of electricity against a \$150
service; the median roof needs \textbf{\$2.44}/kWh to break even. Annual loss cancels
between the value of a clean and the recovery fraction's denominator, so the dollar figure
is identical under all three assumptions while the fraction spans 0.031 to 0.217. What
generalises is the fold geometry, the label-unit diagnosis and the arithmetic.
\end{abstract}
\end{minipage}
\end{center}

\vfill
\clearpage
\twocolumn

% ---------------------------------------------------------------------------
\section{Introduction}
\label{sec:intro}

Global photovoltaic capacity reached roughly \SI{2383}{\giga\watt} by the end of 2025
after a record \SI{510}{\giga\watt} single-year addition \cite{irena2026}, and solar now
supplies more than 8\% of world electricity \cite{iea2026}. US small-scale systems were
projected to generate about 83\,TWh in 2024 \cite{eia2023}, and the
country passed six million individual installations in early 2026 \cite{seia2026}.
Extracting output from arrays already installed is now worth as much as installing more.

\paragraph{Soiling, and why averages mislead.}
Soiling is the accumulation of dust, pollen, and biological film on module glass.
Averaged over fleets the effect is modest: \citet{mejia2013} measured roughly
0.051\%/day across California sites, about 7.4\% by the end of a 145-day dry summer, and
concluded that cleaning a typical residential system does not pay. The distribution
matters more than the mean. In the same study about a quarter of sites soiled at more
than twice the average, and arrays tilted below five degrees soiled roughly five times
faster. Those are the systems a cleaning programme is for, and finding them is a
targeting problem.

\paragraph{The data that would answer it is not available.}
Whether a specific array is soiled is observable from its production history, which
monitored inverters record and vendors retain. A production dip is not yet a cleaning
decision: it does not attribute the loss to dirt, price it over a year, or weigh it
against the cost of a wash. Homes without monitoring have less. Anyone outside the
household, a researcher, a municipality, a cleaning service, has nothing. Today no
residential array is flagged for soiling, clean or filthy.

\paragraph{What we found, stated before the evidence.}
We built the system and it locates arrays well. It cannot rank them, and we initially
believed it could, because our own validation leaked: the fold that produced our reported
out-of-year AUC of 0.710 leaves the held-out year's sites in training through their other
years for 88.7\% of rows. Closing that leak takes the number to 0.622, and latitude and
longitude on their own score 0.644---as well as all forty features. The cause is not the
model. Soiling labels of the kind that exist publicly are measured at monitoring
\emph{stations}, so every rooftop in a station's catchment inherits one label, and no
per-roof feature has anything to fit against.

The obvious repair is per-system labels, and we test it: on telemetry from 735 individual
systems, per-array features do rank systems that share weather, at 0.573 against the 0.5 a
station-label model cannot exceed by construction. But three standard ways of computing that label
put the same systems' median annual loss at 2.72, 9.80 and \SI{20.93}{\percent} and rank
them at Spearman 0.241 to 0.894, and our 0.573 rests on the variant that is in the
minority on ranking. The repair is necessary and not yet sufficient.

Then the result that reorders the others. We measured what a cleaning is worth on 149
metered rooftops, and one wash on the median system recovers \textbf{\$28.10} of
electricity at the retail offset most of them earn, or \$10.14 at the marginal export
rate, against a \$150 service. The median roof would need electricity at
\textbf{\$2.44}/kWh to break even, about five times California retail. That is not a
targeting failure and no better model repairs it. It is the more useful finding, and we
put it last only because it takes the most machinery to establish.

\paragraph{Contributions.}
We ask what a system built \emph{entirely} on public data can and cannot establish.

\begin{enumerate}\itemsep2pt
\item \textbf{A gated detection stage} (\S\ref{sec:methods-inputs}; full protocol in~\cite{gordon2026detection}) on public county
orthoimagery at \SI{0.204}{\metre} ground sample distance, using only permissively
licensed models, with an executable gate definition, confidence frozen on validation,
tile-level bootstrap intervals, and an independent permit-based recall probe
(\S\ref{sec:results-inputs}).
\item \textbf{Per-array roof geometry from public lidar} (\S\ref{sec:methods-inputs}),
validated against installer-reported tilt from utility interconnection filings, plus a
projection trap that silently flattens tilt by \SI{18.9}{\percent} on web-Mercator
point clouds. The
closest prior system, 3D-PV-Locator~\cite{mayer2022pvlocator}, validates tilt per system
against ground truth; ours compares distributions, which is the weaker claim and we mark
it as such.
\item \textbf{A measured account of why station-trained soiling models do not transfer to
individual roofs} (\S\ref{sec:ranking}): a leaking-fold disclosure and its correction, a
variance decomposition, a controlled feature-backfill experiment, a regional holdout, and
a sign reversal with an identified mechanism---together with a positive control showing
that the limit is the label unit rather than the model class.
\item \textbf{A grounded cleaning-economics verdict} (\S\ref{sec:results-econ}): zero of
1{,}865 sites clear, with the one unsourced constant that had been carrying the opposite
conclusion identified and replaced by a directly measured quantity: 505 observed cleaning
events on 149 metered rooftops recover a median \$28.10 of electricity against a \$150
service. We show the value of a clean is invariant to the cleaning assumption while the
recovery \emph{fraction} is not, and report the fraction only as a bracket.
\item \textbf{A reproducibility posture}: every number in this paper names the artifact
that produced it and the command that regenerates it (\S\ref{sec:repro}), and every
constant we could not source is listed rather than buried (\S\ref{sec:unsourced}).
\end{enumerate}

\begin{figure*}[t]
  \centering
  \includegraphics[width=0.98\textwidth]{fig_argument.pdf}
  \caption{The whole paper. Four questions a public-data cleaning product must answer in
  order, what each returned, and the point at which the answer stops mattering. Question~4
  is arithmetic and, \emph{in this market}, it settles questions~2 and~3: a perfect ranker
  cannot create value that is not there. Where the boundary of \S\ref{sec:results-econ}
  is cleared---drier, pricier, slower-soiling---questions~2 and~3 become live again.}
  \label{fig:argument}
\end{figure*}

Figure~\ref{fig:argument} is the argument in one picture, and it is worth stating the shape
before the evidence. Three of the four answers are negative, but they are not
interchangeable. Question~2 fails because of a property of the labels, question~3 fails
because of a property of the measurement, and question~4 fails for a reason that has
nothing to do with either: the quantity of money involved is smaller than the cost of
acting on it. A reader who takes only one thing from this paper should take the last one,
because it is the only one that would not be fixed by a better model.

Figure~\ref{fig:pipeline} shows the pipeline end to end.

\begin{figure*}[t]
  \centering
  \includegraphics[width=0.98\textwidth]{fig_pipeline.pdf}
  \caption{The pipeline. Every input is public and free: county orthoimagery, USGS 3DEP
  lidar, Open-Meteo weather, ESA WorldCover, OpenStreetMap, county parcel
  and permit records, utility interconnection filings, and NREL soiling labels. Both
  vision models are Apache-2.0. Stage boundaries are where the evaluation changes:
  detection is scored on boxes, masks on area, risk on AUC and calibration, and economics
  on expected net dollars.}
  \label{fig:pipeline}
\end{figure*}

% ---------------------------------------------------------------------------
\section{Related Work}
\label{sec:related}

\paragraph{Locating arrays from overhead imagery.}
Public rooftop-PV detection has an established lineage: the Duke
dataset~\cite{bradbury2016duke} released \SI{0.3}{\metre} imagery with array annotations,
and DeepSolar~\cite{yu2018deepsolar} built a national deployment database by
classification with weak localization. Our detector consumes that lineage rather than
extending it. It uses RF-DETR~\cite{rfdetr2025}, a real-time DETR
variant~\cite{carion2020detr}, with an optional SAM2~\cite{ravi2024sam2,kirillov2023sam}
mask stage, both Apache-2.0; the licence matters because the widely used YOLO
implementations are AGPL-3.0, which obstructs reuse by exactly the municipalities and
nonprofits a public-data pipeline is for. Two requirements differ from the prior work and
are specified in~\cite{gordon2026detection}: downstream stages consume \emph{area}, so we
need instance polygons rather than image-level labels, and the gate is evaluated at the
tile level through the production inference path rather than on crops. Detection is an
input here, not a contribution.

\paragraph{Per-array geometry.}
3D-PV-Locator~\cite{mayer2022pvlocator} recovers tilt and azimuth per system from 3D
building data and validates them per system against ground truth. We fit planes to public
lidar returns inside each array polygon and can only validate distributionally, because the
interconnection filings that report tilt carry no parcel, address or coordinate. We say so
where the comparison appears rather than presenting the weaker design as equivalent.

\paragraph{Soiling measurement and modelling.}
Soiling loss is extracted from yield directly by \citet{deceglie2018}, and NREL publishes a
station-level US soiling map~\cite{nrel_soiling_map,micheli2021} reporting an
insolation-weighted soiling ratio (IWSR). Empirical models include
Kimber~\cite{kimber2006}, which accumulates loss with particulate concentration and resets
on rain, and SOMOSclean~\cite{micheli2021}, an exponential approach to a saturation
ceiling with rain-proportional cleaning. We use both as \emph{features} under measured
IWSR labels, and report in \S\ref{sec:risk} why the converse configuration is a leakage
trap we fell into and discarded. \citet{porcar2018} confirmed biofilm colonisation on
panels in Berkeley but measured no power loss; that gap is the open channel of
\S\ref{sec:openchannel}.

\paragraph{Uncertainty.}
Our loss intervals are conformalised quantile regression~\cite{romano2019cqr} over
gradient-boosted trees~\cite{chen2016xgboost}, because the raw quantile heads
under-covered badly. Degradation and soiling-ratio decomposition use
RdTools~\cite{rdtools} on PVDAQ~\cite{pvdaq}; cell-level occlusion physics uses the
Bishop diode model~\cite{bishop1988} via pvlib~\cite{holmgren2018pvlib}.

% ---------------------------------------------------------------------------
\section{Methods}
\label{sec:methods}

\subsection{Inputs: array detection and roof geometry}
\label{sec:methods-inputs}

The pipeline's first two stages supply the units this paper analyses: array
polygons and their tilt. Both are specified in full in~\cite{gordon2026detection}. Every
number this paper relies on is stated here, so nothing below depends on reading that
report; it is cited for protocol detail, not for evidence.

\paragraph{Detection.} RF-DETR at 728-pixel input on \SI{21}{\centi\metre} county
orthoimagery, with an optional SAM2 mask stage, both permissively licensed. A full
tile is cut into a $2\times2$ grid of 640-pixel chips, each chip gets one
model pass, and detections are merged across the seams by NMS at IoU 0.55. The
acceptance gate is stated as an executable definition rather than a number---the
script is the specification---and requires tile-level box-$F_1 \geq 0.75$, a
bootstrap CI lower bound $\geq 0.70$, and recall $\geq 0.70$, with the confidence
threshold tuned on validation and frozen before test is scored.

\paragraph{Roof geometry.} Per-array tilt and azimuth are fitted to \SI{1}{\metre}
public 3DEP lidar returns inside each array polygon. One detail is load-bearing and
easy to get wrong: fitting a plane to point clouds reprojected to web-Mercator
reports tilt \SI{18.9}{\percent} \emph{too shallow}, because the projection stretches
horizontal distance against an unstretched vertical axis, shrinking $dz/dx$. Fits are performed in a local metric CRS.
Arrays without a usable fit keep a flagged null rather than an imputed angle.

\subsection{Soiling risk model}
\label{sec:risk}

\paragraph{Labels.}
We train against the NREL PV Soiling Map~\cite{nrel_soiling_map}, whose IWSR is the
fraction of expected insolation-weighted energy actually delivered, so \emph{lower is
dirtier}. The annual panel release gives \textbf{891 station-year rows across 146 stations
spanning 2008--2022}. A further \textbf{111 stations report only a censored summary}
(IWSR $>0.99$, known-clean without an annual breakdown); we retain these at the censoring
value at half the weight of a panel row, giving \textbf{1{,}002 rows across 257 stations}.
Rows are binarized at IWSR $<0.97 \rightarrow$ at-risk, yielding a 49.4\% base rate.
Panel rows span six states (AZ, CA, GA, MD, NC, OR); the full station set including
summary rows spans 15. The two counts are not interchangeable.

\paragraph{The central methodological compromise.}
We train on instrumented ground stations and apply the result to residential rooftops.
That transfer is unvalidated. It is also the only option available without production
data, and \S\ref{sec:ranking} measures how badly it binds.

\paragraph{Features.}
Thirty-four features in three families, from public sources: \emph{weather} from
Open-Meteo~\cite{openmeteo} over trailing 7/30/90-day windows (precipitation, maximum
\SI{10}{\metre} wind, mean relative humidity, maximum temperature) plus two rain-timing
features, consecutive dry-day streak and days since the last \SI{5}{\milli\metre} rain
event; \emph{location and structure}, latitude, longitude, month, tilt, elevation,
land-cover class, distance to nearest highway and to nearest agricultural parcel, and five
ESA WorldCover~\cite{zanaga2022worldcover} fractional covers; and \emph{physics proxies}
from Kimber~\cite{kimber2006} and SOMOSclean~\cite{micheli2021}, each as an instantaneous
value plus 7/30/90-day means, with SOMOSclean's saturation and rate re-fit to coastal
California stations ($s_{\text{sat}}=0.08$, $k=15$).

Air quality was tested and dropped. PM$_{2.5}$ and PM$_{10}$ rolling means covered 21\% of
label rows (208 of 1{,}002) and received zero splits in the fitted model, so the family
contributed no gain at any coverage level and is excluded from the feature count above.

\paragraph{A leakage trap.}
Physics proxies as features under measured IWSR labels are leakage-free. The converse is
not. Ablations that used SOMOSclean-\emph{derived} labels while retaining the SOMOSclean
feature returned cross-validated AUC of exactly 1.0: the feature was the label. Those runs
are discarded.

\paragraph{Model and validation.}
XGBoost~\cite{chen2016xgboost}: 200 estimators, max depth 4, learning rate 0.05, subsample
0.8, column subsample 0.8, $L_2$ regularization 5.0. Depth and $L_2$ were tuned against
overfitting, because this many features on $\sim$900 rows overfit visibly at depth 5. Outputs pass
through an isotonic calibrator so the surfaced number reads as a probability.
\emph{Spatial} validation clusters stations within \SI{10}{\kilo\metre} into the same
fold. \emph{Temporal} validation holds out each of the 15 panel years in turn and pools
every row's out-of-its-own-year prediction; \S\ref{sec:results-risk} explains why this
replaced a single-year holdout. \emph{Calibration} is evaluated out-of-year by fitting
each training year's isotonic map and applying it to the held-out year.

\subsection{Loss magnitude and cleaning economics}
\label{sec:econ}

\paragraph{From probability to magnitude.}
Multiplying a calibrated classification probability by a constant to obtain a loss
percentage is a category error rather than a bad estimate. We instead use
three XGBoost quantile heads (0.10/0.50/0.90) trained on $(1-\text{IWSR})\times100$ under
the same spatial CV and conformalised~\cite{romano2019cqr}. No fallback path from risk
score to loss percentage exists any more, deliberately; callers without a magnitude fall
back to a flagged base rate.

\paragraph{The decision model.}
Recovered energy is priced at the site's modelled marginal tariff and its system size, against the cost
of the chosen cleaning method, over 2{,}000 Monte Carlo draws per site propagating the
loss prediction interval, an area uncertainty ($1\sigma = 0.26$, the measured
21\,cm/60\,cm area shift), the rate band, and recovery timing. Economics run \emph{per
site}, not per detected polygon: \SI{21}{\centi\metre} imagery fragments a roof into 1.80
polygons per site against 1.30 at \SI{60}{\centi\metre}, and per-polygon costing
overcharged a four-way-split roof by $2.3\times$.

\paragraph{The constant that decided everything.}
\texttt{recovery\_frac}, the share of a year's soiling loss that one cleaning recovers,
had been \emph{assumed} at 0.90 and never estimated. We estimate it by simulating the
soiling trajectory with and without a wash over the site's actual rain record
(\S\ref{sec:results-econ}).

\paragraph{Tariffs.}
The rate had been a flat \$0.25/kWh. We replace it with a
bill-reconciled PG\&E E-TOU-C plus community-choice stack, validated to
$\pm$\$0.22/month over 11 consecutive real bills, giving a production-weighted
retail offset of \$0.4573/kWh and a net-billing export credit of
\$0.0392/kWh~\cite{cpuc_nbt}. Only 5.5\% of a south-facing array's
annual output lands in the 4--9\,pm peak, so roughly 94\% of soiling-lost kWh are
off-peak. Site tariff vintage is joined from public interconnection
filings~\cite{dgstats}.

% ---------------------------------------------------------------------------
\section{Results}
\label{sec:results}

\subsection{Inputs, in brief}
\label{sec:results-inputs}

The detection and roof-geometry stages are reported in full in the companion
paper; what matters here is that the inputs to the soiling analysis are gated
rather than assumed. Detection reaches tile-level box-$F_1$ \textbf{0.826}
(95\% CI [0.798, 0.853], precision 0.850, recall 0.803, $n_{\text{gt}}=585$) at a
confidence tuned on validation and frozen before test was scored, using only
permissively licensed models. Recall measured independently against county
building permits the detector never saw is \textbf{73.6\%} on the imaged era,
which is the number that bounds the funnel: a missed array is a missed roof, and
no downstream model recovers it. The mask stage contributes area rather than
detection, and its one dangerous failure---segmenting the whole roof instead of
the array---is a prompt-box problem, eliminated at the production 15\% box shrink
(median mask IoU 0.844, median area/GT 1.01, roof-grab 0.0\%). Per-array tilt comes
from public lidar. Its validation is distributional rather than paired---the
interconnection records that report tilt carry no parcel, address or coordinate,
so no array can be matched to the filing describing it. The medians agree to
\SI{0.03}{\degree} ($n=3{,}068$ lidar fits against $8{,}573$ filings) while a
Kolmogorov--Smirnov test rejects the distributional match at $p<10^{-90}$: every
reported tilt is an integer, five values cover 57.6\% of filings, and 27\% are
exactly \SI{18}{\degree}, the 4:12 roof pitch. Installers file a nominal pitch off
a short menu; lidar measures a continuous angle.
That limitation is one of three reasons tilt is deliberately excluded from the
risk model in \S\ref{sec:ranking}.

\subsection{Soiling risk: the gate it cleared was leaking}
\label{sec:results-risk}

\paragraph{Which model these numbers describe.}
Two feature sets appear in this work and it matters which is which. The \emph{published}
configuration carries 40 features including six air-quality columns; the \emph{honest}
production set is the same 40 minus those six, because they cover 20.8\% of rows and
contribute nothing measurable. Every corrected result below is reported for both, and they
are indistinguishable: on the joint fold, 0.6222 for 40 features and 0.6228 for 34. We
quote the 40-feature figure throughout, so that the comparison against our own published
numbers is like-for-like, and give the 34-feature value wherever it differs. Where this
paper says ``40 features'' it means the published configuration, not a recommendation.

\paragraph{Spatial and temporal validation.}
Spatially cross-validated AUC is \textbf{0.718} ($0.719 \pm 0.003$ over seven seeds) for
the 34-feature set under the current configuration. The 40-feature set that also carried
air quality reads 0.712 ($0.721 \pm 0.006$), so dropping that family moves the metric less
than the seed spread. An earlier stored figure of 0.728 for the 40-feature set no longer
reproduces and should not be quoted. Against the figure we report, the margin over our
0.70 acceptance bar is under two points, and against the 40-feature set it is one.

Our initial temporal design held out 2022 as a single test year and returned AUC 0.679 on
97 rows. That estimate is unusable, and we report it because it is the kind of number
over-read in both directions: its Hanley--McNeil standard error is 0.061, its 95\% CI is
[0.557, 0.797], and the 0.70 bar sits 0.32 standard errors away. A 97-row year cannot
distinguish 0.68 from 0.70. Replacing it with rolling leave-one-year-out pooling over
$n=891$ (440 at-risk, 451 not) gives \textbf{AUC 0.710}, SE 0.017, bootstrap 95\% CI
\textbf{[0.676, 0.742]}, $P(\text{AUC}\geq0.70)=0.70$: a $3.5\times$ tighter estimator on
identical data. \textbf{That estimator is also leaking, and the leak is the finding.}
Leave-one-year-out holds out a year but keeps the held-out year's sites in training
through their other years: measured, \textbf{88.7\%} of held-out rows have their own site
in the training set. The 10\,km spatial fold has the mirror defect, grouping by location
while mixing years. Neither is a test of a roof we have not seen in a year we have not
seen. Building that fold---hold out year $y$, then drop from training every row whose
site appears in $y$---gives \textbf{AUC 0.622}, 95\% CI \textbf{[0.571, 0.670]},
$P(\text{AUC}\geq0.70)=\textbf{0.001}$, on the same 891-row pooled basis. \emph{The model
does not clear its temporal gate.} Latitude and longitude alone score 0.644 on that fold,
at or above the full 40-feature set, which is the identifiability problem appearing
directly in the metric: station-level labels give every roof in a catchment one label, so
no per-roof feature has anything to fit.

\paragraph{What the prior literature gets, scored the same way.}
The natural next question is what a published soiling model achieves on this task. Two are
already computed over our rows---Kimber~\cite{kimber2006} and
SOMOSclean~\cite{micheli2021}. Kimber is used exactly as published, with no parameter
estimated on our data, so it needs no fold: there is nothing fitted to leak. SOMOSclean is
\emph{not} in that position, and we flag it rather than lean on it---its saturation and
rate were re-fit to coastal-California stations (\S\ref{sec:risk}), some of which are rows
scored here, so its 0.544 should be read as an optimistic bound on an unfitted SOMOSclean
rather than as a clean baseline. Kimber carries the comparison. Scored on the same 891 panel rows,
Kimber's instantaneous IWSR proxy gives \textbf{AUC 0.610} (95\% CI [0.566, 0.652]) and
SOMOSclean's gives 0.544. Handing all seven physics columns to the learner under the joint
fold gives 0.5706.

Set that beside our own numbers on the identical basis---0.622 for forty engineered
features, 0.644 for latitude and longitude---and the ordering is the result. A 2006
empirical formula with no fitting is within noise of a gradient-boosted model trained on
fifteen panel years of weather, air-quality and static features, and two coordinates match
or beat both. This is
the same conclusion the feature ablations reach, arrived at from outside: on
station-level labels there is very little to learn beyond where the station is.

\paragraph{Is the drop the leak, or the smaller training set?}
Closing the leak also deletes data: holding out year $y$ and then removing every row whose
site appears in $y$ costs a median \SI{28.8}{\percent} of the leave-one-year-out training
set, and \SI{74.3}{\percent} in the worst year. Some of the fall could therefore be volume
rather than leakage, and the control is cheap. For each year we build a third training set
of the \emph{same size} as the joint fold's that \emph{keeps} the leak: start from the
leave-one-year-out set and delete an equal number of rows drawn only from sites absent from
year $y$. Such deletions cannot remove a leaking row, because the leaking rows are exactly
those from sites present in $y$.

\begin{center}\small
\begin{tabular}{@{}lr@{}}
\toprule
Training set & pooled AUC \\
\midrule
leave-one-year-out, leak present, full data & 0.7144 \\
size-matched control, leak present          & 0.6986 \\
joint out-of-site and out-of-year           & \textbf{0.6222} \\
\bottomrule
\end{tabular}
\end{center}

\noindent
Of the 0.0922 total, \textbf{0.0764 (\SI{83}{\percent}) is attributable to the leak and
0.0158 (\SI{17}{\percent}) to training size}. Bootstrapping the decomposition over stations
puts the leak share at 95\% CI \textbf{[52, 99]\%} and the leak's own contribution at
$[+0.031, +0.121]$, which excludes zero. The split is therefore not precisely resolved at
this sample size: what the data support is that the \emph{majority} of the fall is
leakage, not that the share is 83\%. The five-seed spread on the control itself is 0.0012,
so the imprecision is sampling variance in the AUCs rather than which rows were deleted.

That 0.622 is not a fragile number. Across seven seeds it is $0.6201 \pm 0.0047$;
penalised logistic regression, a different model class entirely, gives 0.5707, which is
lower rather than higher; and the weighting of the censored summary rows---the choice a
referee is most likely to query---moves it only between 0.594 (dropped) and 0.622
(production). No setting we tried approaches the 0.70 bar. The regional result is equally
insensitive to fold geometry, ranging 0.638 to 0.664 as the number of k-means regions
varies from 5 to 12.

We report the leaking numbers alongside the corrected ones because reproducing them
requires knowing which of three reporting bases is in use, and none was previously
written down: pooled against mean-of-folds (about 1 point), all 1{,}002 rows against the
891 panel rows (3 to 6 points, since the 111 censored summary rows are easier to classify
and all carry a single sentinel year), and a plain gradient-boosted fit against the
calibrated training pipeline (about 0.4 points). Matched on basis, every previously
published figure reproduces to within 0.012. Out-of-year calibration is retained:
Brier 0.218 against a base-rate reference of 0.250, ECE 0.046, MCE 0.126
(Figure~\ref{fig:risk}). The weakest year is 2019 (AUC 0.626); 2008--2012 are near
single-class.

\begin{figure*}[t]
  \centering
  \includegraphics[width=0.98\textwidth]{fig_risk_validation.pdf}
  \caption{Out-of-year reliability, and the two folds side by side. \emph{Left:}
  calibration survives temporal shift. \emph{Right:} the published leave-one-year-out
  interval, whose fold leaves the held-out year's sites in training for \SI{88.7}{\percent}
  of rows, against the same model scored jointly out-of-site and out-of-year. The published
  interval straddles the acceptance bar; the corrected one sits below it, with
  $P(\text{AUC}\geq0.70)=0.001$.}
  \label{fig:risk}
\end{figure*}

\subsection{Why the model cannot rank individual roofs}
\label{sec:ranking}

This is our second negative result, and it reframes what the pipeline claims. The model is
valid for \emph{level} and not for \emph{ranking homes against each other}.

\paragraph{Variance decomposition.}
Across the 1{,}710 residential-scale sites (2--\SI{15}{\kilo\watt}), the per-home dollar
figure is driven by system size (p10--p90 ratio $3.20\times$, \textbf{86.7\%} of variance)
and roof orientation ($1.30\times$, \textbf{11.1\%}); the soiling model contributes
$1.15\times$ and \textbf{2.2\%} (Figure~\ref{fig:variance}). \textbf{97.8\% of the spread
comes from what the detector and the lidar measure.} Doubling the soiling model's spread
would add less differentiation than joining one public permit file. The decomposition
mixes per-array and per-site quantities, so it estimates rather than partitions; the ratio
is scale-invariant and therefore unaffected by the level calibration below.

\begin{figure*}[t]
  \centering
  \includegraphics[width=0.98\textwidth]{fig_variance_decomposition.pdf}
  \caption{Left: share of per-home dollar variance by driver. Right: the soiling head's
  entire within-AOI spread, 0.38 points p10--p90 across 1{,}710 homes.}
  \label{fig:variance}
\end{figure*}

\paragraph{The mechanism.}
NREL labels are station-level. Across a \SI{20}{\kilo\metre} coastal AOI, feature standard
deviations are 0.3\% to 23\% of the training standard deviations, and \textbf{58.1\% of
model importance sits on features effectively constant across the AOI}. A model trained to
separate 15 states has almost nothing left when shown 1{,}865 homes in one town, because
the training set contains \emph{zero} examples of per-roof variation.

\paragraph{The controlled experiment.}
Several features were missing from the deployment matrix and silently median-filled, an
obvious candidate explanation with an obvious fix. We supplied the WorldCover features
properly, resolving all 3{,}349 points. The within-AOI p10--p90 spread moved from 0.776 to
0.758 points, which is to say not at all, and dead importance moved 64.9\% to 58.1\%. The
plumbing gap was not the cause.

\paragraph{The sign reversal.}
Arrays on tree pixels scored \emph{lower} (5.36\%) than built-up (5.55\%). The model
learned \texttt{worldcover\_tree} from NREL stations, where a tree pixel means ``forested
rural region, wetter air, less dust,'' not ``a tree overhanging this roof.'' \textbf{Same
feature name, different referent.} We predicted and then confirmed the same failure for
tilt, and deliberately excluded our measured per-array tilt from the risk model for three
measured reasons: the training interquartile range $[\SI{7.49}{\degree},
\SI{12.00}{\degree}]$ contains only 8.9\% of our roofs; gradient-boosted trees do not
extrapolate, so the deepest tilt split is at \SI{26.0}{\degree} and predictions at
\SI{30}{\degree}, \SI{35}{\degree} and \SI{40}{\degree} are identical; and the training
feature is effectively a station identifier (42.4\% of rows at exactly
\SI{7.49}{\degree}, 24 unique values, 0.372 correlation with longitude). The learned
effect accordingly carries the \emph{wrong sign}. Tilt belongs in the physics, where a
\SI{30}{\degree} roof self-cleans on less rain than a \SI{5}{\degree} one, not in a model
that never saw per-roof tilt vary.

\begin{figure*}[t]
  \centering
  \includegraphics[width=0.98\textwidth]{fig_label_unit.pdf}
  \caption{The identifiability problem, measured rather than argued. \emph{Left:} every
  metered residential system inside one \SI{0.5}{\degree} cell. Their annual soiling losses
  span 0.23 to \SI{5.35}{\percent}, and a station-level label gives all 103 of them the
  single value marked in red. \emph{Right:} the cell is not unusual. Across all 37 cells
  holding five or more metered systems, the within-cell standard deviation is a median
  \SI{42}{\percent} of the cell mean. Roughly half the variation a per-roof model would
  need to explain is variation that station labels discard by construction.}
  \label{fig:labelunit}
\end{figure*}

\paragraph{The mechanism, measured on metered roofs.}
Figure~\ref{fig:labelunit} shows what the label unit costs. One \SI{0.5}{\degree} cell in
our fleet holds 103 metered residential systems whose measured annual soiling losses range
from 0.23 to \SI{5.35}{\percent}, a factor of 23. A station-level label assigns all 103 the
same number. That cell is not a special case: across the 37 cells holding at least five
metered systems, the within-cell standard deviation is a median \SI{42}{\percent} of the
cell mean. Approximately half of the variation a per-roof model would have to explain is
removed before training begins, and no feature engineering reaches it.

\paragraph{What per-system labels recover.}
Everything above is a diagnosis without a remedy, and a reader is entitled to ask whether
\emph{any} public-data model can rank roofs or whether the task is impossible. It is not.
The decisive test is ranking systems that share weather: two arrays in the same
\SI{0.5}{\degree} cell see the same climate, air quality and rainfall, so a station-label
model assigns them identical features and emits identical predictions. Its within-cell
pairwise concordance is therefore exactly 0.5 \emph{by construction}, not as an empirical
finding. Anything above 0.5 must come from per-array information.

We test this on the PVDAQ fleet, where the labels are per system rather than per station:
94 cells holding two or more systems, 735 systems, 11{,}653 same-cell pairs, evaluated
leave-one-cell-out so the model never sees the cell it ranks (Figure~\ref{fig:persystem}).
Per-array features alone---tilt, azimuth, capacity, age---give concordance \textbf{0.573}
(95\% CI [0.512, 0.609], $P(>0.5)=0.995$, bootstrapped over cells because systems within a
cell are correlated). Location alone, the station-model analogue, gives \textbf{0.484},
indistinguishable from chance even with sub-cell coordinates available. Per-array
information is therefore present and learnable, and the station-label ceiling is
structural rather than incidental.

\begin{figure}[t]
  \centering
  \includegraphics[width=0.96\columnwidth]{fig_per_system_ranking.pdf}
  \caption{Within-cell pairwise concordance on systems that share weather. A station-label
  model cannot exceed 0.5 here by construction, and per-array features clear it on
  \texttt{perfect\_clean} labels. The result does not survive the alternative cleaning
  assumption (\texttt{half\_norm\_clean}, concordance 0.512), which is why we report the
  label instability rather than the ranking as the finding.}
  \label{fig:persystem}
\end{figure}

\paragraph{But the per-system label is not yet well posed, and this is the more useful
finding.} The result above uses the \texttt{perfect\_clean} soiling ratio. RdTools offers
two further cleaning assumptions, and on 146 California residential systems where all three
fits are non-degenerate they disagree about both the level and the ordering:

\begin{center}\small
\begin{tabular}{@{}lrrr@{}}
\toprule
assumption & median loss & IQR & median CI width \\
\midrule
\texttt{perfect\_clean}   & \SI{2.72}{\percent} & [1.95, 3.74]  & 0.63 \\
\texttt{half\_norm\_clean} & \SI{9.80}{\percent} & [8.53, 11.98] & 6.92 \\
\texttt{random\_clean}    & \SI{20.93}{\percent} & [17.34, 24.33] & 19.84 \\
\bottomrule
\end{tabular}
\end{center}

\noindent
Spearman between them, on these 146 systems with all three fits non-degenerate and a
tighter Monte Carlo: \texttt{perfect}--\texttt{half\_norm} \textbf{0.241},
\texttt{perfect}--\texttt{random} \textbf{0.318}, and
\texttt{half\_norm}--\texttt{random} \textbf{0.894}. Re-running the within-cell experiment
against \texttt{half\_norm\_clean} gives concordance \textbf{0.512}, at the edge of chance.
(The \texttt{perfect}--\texttt{half\_norm} correlation is 0.126 across the larger
735-system ranking sample and 0.241 here; the two populations and Monte Carlo budgets
differ, and we quote 0.241 throughout because it is the run in which all three assumptions
were computed on identical systems.)

\emph{This is not one bad estimator.} That was our first reading and it does not survive
the third assumption: \texttt{half\_norm} and \texttt{random} are different assumptions and
they rank the same roofs at 0.894, while \texttt{perfect\_clean} disagrees with both.
The uncomfortable part is that the plausible \emph{level} and the majority \emph{ranking}
point at different estimators. Only \texttt{perfect\_clean} lands near the independent
references for this population---\SI{2.76}{\percent} for PVDAQ systems near the pilot,
\SI{4.70}{\percent} for NREL's coastal-California stations---and that comparison is partly
circular, because NREL's IWSR is itself defined under a perfect-cleaning assumption. Our
own 0.573 therefore rides on the minority ranker, and we say so rather than let a reader
discover it.

The honest statement is narrower than the section began. Per-array signal exists, per-array
covariates exist, and station labels can reach neither. But a per-system label recovered
this way is \emph{necessary and not yet sufficient}: the choice of cleaning assumption
moves the level eightfold and can move the ranking to chance, and no independent reference
can adjudicate it, because the references available share an assumption with one of the
candidates. Closing that gap---an estimator both unbiased for recoverable loss and tight
enough to order two roofs in one neighbourhood---is the prerequisite for per-roof soiling
targeting. It is a measurement problem, and it is the one worth working on. We flag it
because the obvious next move in this literature is to scale per-system label extraction,
and scaling a label whose standard variants disagree this much will produce confident and
unreproducible rankings.

\begin{figure*}[t]
  \centering
  \includegraphics[width=0.98\textwidth]{fig_coverage.pdf}
  \caption{Why no model here should be expected to travel. The 891 annual panel rows that
  carry temporal structure are \SI{89.6}{\percent} west of $-\SI{114}{\degree}$, and
  \SI{91}{\percent} of the 146 stations behind them. The eastern United States appears in
  this dataset almost entirely as \emph{censored summary rows}---one row per station, no
  annual breakdown, \SI{87}{\percent} of them eastern---which contribute nothing to a fold
  that holds out a year. Every metered rooftop in our economics is Californian.}
  \label{fig:coverage}
\end{figure*}

\paragraph{The dataset is less geographically balanced than a station count suggests.}
Counted over all 257 stations, \SI{57}{\percent} sit west of $-\SI{114}{\degree}$, which
sounds tolerable. Counted over the 146 stations that contribute \emph{panel} rows---the
only ones with multiple years, and therefore the only ones a temporal fold can use---the
figure is \SI{91}{\percent} (Figure~\ref{fig:coverage}). The eastern representation is
almost entirely the 111 censored summary rows, \SI{87}{\percent} of which are eastern and
each of which contributes a single observation with no annual structure. A geographic
holdout on this dataset is therefore closer to a within-Southwest holdout than the station
count implies, and we read our own out-of-region numbers accordingly.

Two caveats on those two numbers, both of which we found by
re-running the experiment. The fold holds out whole sites, since a region contains them
entirely, but it does not hold out years; and the script read its hyperparameters from a
configuration key that does not exist, so it had always fitted a stock model rather than
the one we ship. Correcting both gives pooled \textbf{0.655} and \textbf{0.527} on the
largest region. The direction of the conclusion is unchanged and its magnitude is worse. \SI{89.6}{\percent} of the panel rows sit west of $-\SI{114}{\degree}$, so the model has
effectively already seen the Southwest. This model
does not travel to an unseen region. More features will not fix it; a geographically
balanced label set will.

\begin{figure*}[t]
  \centering
  \includegraphics[width=0.98\textwidth]{fig_regional_holdout.pdf}
  \caption{Out-of-region generalization. The largest held-out region scores near chance,
  and no feature set recovers it.}
  \label{fig:regional}
\end{figure*}

\paragraph{Level is biased high, and we correct the symptom, not the cause.}
The raw head predicts 5.53\% annual recoverable loss for this AOI. PVDAQ residential
systems within \SI{120}{\kilo\metre} measure a median of 2.76\% ($n=118$). We apply a
multiplicative calibration of 0.500 to all three quantiles, placing the AOI at 2.76\%
against NREL's own coastal-California station median of 2.80\%, an independent reference
not used to derive the factor. The cause is the training geography: 774 of 891 rows are
California, dominated by Los Angeles and San Bernardino, which is desert. The fix is a
per-array label set, not a scalar.

\subsection{The cleaning-value verdict}
\label{sec:results-econ}

\textbf{Zero of 1{,}865 detected Santa Cruz sites show a positive expected net from a
cleaning}, on recoverable soiling, at any electricity rate to
\$0.70/kWh, at any system size, over 2{,}000 Monte Carlo draws each.
Maximum $P(\text{net}>0)$ across all sites is 0.0000. Figure~\ref{fig:shortfall} shows the shortfall, cleaning cost divided by annual recovered
dollars, for every site.

\paragraph{The same verdict on measured rooftops, with nothing assumed.}
The result above rests on a modelled recovery fraction and a detected, not metered,
array. We can now remove both. For \textbf{149} residential California systems in PVDAQ
we measure recovery directly from RdTools soiling intervals---505 cleaning events falling
in dry windows, separated from rain resets by daily precipitation---together with each
system's own annual soiling loss, its observed heavy-rain frequency, and its nameplate
capacity. \emph{What we measure is the value of a clean, not a recovery fraction.} One
wash recovers a median \textbf{241 point-days} of performance (IQR [160, 340]), worth
\textbf{\$28.10} at the retail offset and \$10.14 at the export rate (IQR [17.32, 43.70]
and [6.25, 15.77]), against a median \$150 service. Those are assumption-free: they depend
on the observed step and its duration, not on any model of unobserved cleaning.

\paragraph{The dollar verdict does not depend on any of that.}
\S\ref{sec:ranking} shows the per-system \emph{label} is assumption-dependent by a factor
of eight. It would be reasonable to expect the economics to inherit that, and they do not,
for an algebraic reason worth stating explicitly. Recovered dollars are the annual loss in
dollars multiplied by the recovery fraction, and the recovery fraction carries that same
annual loss in its denominator:
\begin{equation*}
\underbrace{kW \cdot h \cdot \eta \cdot p \cdot \tfrac{L}{100} \cdot 365}_{\text{annual loss in \$}}
\;\times\;
\underbrace{\frac{\Delta \cdot d}{\tfrac{L}{100} \cdot 365}}_{\text{recovery fraction}}
\;=\; kW \cdot h \cdot \eta \cdot p \cdot \Delta \cdot d .
\end{equation*}
The annual loss $L$ cancels exactly. What a clean is worth depends only on the measured
performance step $\Delta$, how long it persists $d$, the system size, the insolation and
the tariff---none of which involves an assumption about unobserved cleaning. Verified on
every system rather than one: across all 149 rooftops and all three cleaning assumptions,
spanning an eightfold range of assumed annual loss, the largest deviation in recovered
dollars is $4.3\times10^{-14}$ dollars, which is floating-point noise. \emph{The contribution most exposed to the label problem is the
ranking; the one that decides the product is immune to it.}

\paragraph{Recovery fraction is a ratio, and we demote it accordingly.}
The literature's constant---assumed at 0.90, modelled here at 0.045---is
$\Delta d / L$, a measured numerator over an \emph{assumption-dependent} denominator. By
the same cancellation that makes the dollars robust, the ratio inherits the full label
instability of \S\ref{sec:ranking}. On the 112 systems carrying all three cleaning
assumptions:

\begin{center}\small
\begin{tabular}{@{}lrr@{}}
\toprule
denominator & median $L$ & implied recovery fraction \\
\midrule
\texttt{perfect\_clean}    & \SI{2.97}{\percent}  & 0.217 [0.132, 0.361] \\
\texttt{half\_norm\_clean} & \SI{9.61}{\percent}  & 0.063 [0.048, 0.097] \\
\texttt{random\_clean}     & \SI{20.58}{\percent} & 0.031 [0.022, 0.044] \\
\bottomrule
\end{tabular}
\end{center}

\noindent
A sevenfold range, entirely from the denominator. We therefore do \emph{not} report a
single measured recovery fraction, and an earlier draft of this paper was wrong to. What
survives is one half of a bracket: the \textbf{0.90 originally assumed is above all three},
so the original business case was wrong by more than the uncertainty in any direction. The
modelled 0.045 is \emph{not} below all three---\texttt{random\_clean} implies 0.031---so it
falls inside the measured range rather than beneath it, and we cannot claim the simulation
is conservative against every denominator. Anyone recomputing from the published fleet label file obtains 0.217, that
file carrying \texttt{perfect\_clean}.
At the marginal export rate of \$0.165/kWh, \textbf{zero of 149} systems clear; at the
full retail offset of \$0.4573/kWh, two do. The median system requires
\textbf{\$2.44/kWh}, roughly five times California retail (Figure~\ref{fig:value}, left).
Peak-sun-hours are measured per site from the cached irradiance record rather than held at
the \SI{5.5}{\hour} constant the pipeline previously applied everywhere; that constant is a
Santa Cruz value, the measured median across 29 cells is \SI{5.09}{\hour}, and the constant
is too high for 22 of them---so it had been making the cleaning case optimistic rather than
conservative.

The obvious objection is our \$150 service floor, which is unsourced and which 135 of the
149 systems price at. It is not load-bearing, because what a wash returns is
too small for any cost structure to reach. One wash on the median system recovers
\textbf{\$28.10} at the retail offset and \$10.14 at the export rate
(Figure~\ref{fig:value}, right). A wash would have to cost under \$10 before half of these
systems cleared at the export rate; at \$25, below any plausible service price, 9.4\%
clear.

\emph{This comparison is per-wash on both sides, which is what makes it robust.} A
homeowner buys one wash at a time, each worth $\Delta d$ priced, and by the cancellation
above that value carries no cleaning assumption. We previously described the same figure as
``the entire annual recoverable value'', which was a mislabel rather than a different
number: the two differ only by using the \SI{5.5}{\hour} constant instead of each site's
measured insolation, and agree to $4\times10^{-14}$ dollars per system when computed
alike. An annual prize \emph{would} depend on the cleaning assumption, and we make no claim
about one.

\paragraph{Where cleaning would pay.}
Stating where the product fails is weaker than stating where it would work, so we sweep
the boundary (Figure~\ref{fig:surface}). Rain frequency enters twice and both times
against the seller: more rain means less accumulated loss to recover, and a shorter
benefit window before the next natural reset. Both quantities come from the same daily
trajectory, so they compound rather than multiplying as independent constants. Reading the boundary off the full grid, a positive expected net requires \emph{all} of a
re-soiling time constant $k \geq \SI{30}{\day}$, at most 15 heavy-rain resets per year, a
tariff above \$0.20/kWh and a system above \SI{4}{\kilo\watt}. Each threshold is the most
permissive value at which any grid cell clears while the other axes are free, so they are
necessary conditions read jointly, not a recipe: only 323 of 49{,}920 cells are positive,
and Figure~\ref{fig:surface} plots one deliberately generous slice of that grid rather
than the thresholds themselves. The surface is
non-monotonic in $k$: too fast and there is no benefit window, too slow and the array
never soils enough to be worth cleaning. Coastal California sits at $k=\SI{15}{\day}$ with
16.5 resets per year, far inside the losing region.

\begin{figure*}[t]
  \centering
  \includegraphics[width=0.98\textwidth]{fig_cleaning_value.pdf}
  \caption{The cleaning-value verdict on 149 measured rooftops, with recovery measured
  rather than modelled. \emph{Left:} the electricity price each system would need for one
  professional clean to break even; the median needs \$2.44/kWh against a \$0.165
  marginal export rate. \emph{Right:} what one wash returns per system at both the export
  rate and the full retail offset, against the cost of the service. Both sides of that
  comparison are per-wash, and the left-hand quantity carries no cleaning assumption. What
  a wash returns is smaller than any plausible price under either tariff.}
  \label{fig:value}
\end{figure*}

\begin{figure}[t]
  \centering
  \includegraphics[width=0.98\columnwidth]{fig_breakeven_surface.pdf}
  \caption{Where one clean would pay, under three concessions stacked in the product's
  favour: \$0.70/kWh, a \SI{20}{\kilo\watt} roof, and a 12-point saturated soiling level.
  The measured operating point remains deep in the losing region, and at the measured
  median system size no grid cell is positive.}
  \label{fig:surface}
\end{figure}

\begin{figure*}[t]
  \centering
  \includegraphics[width=0.98\textwidth]{fig_shortfall.pdf}
  \caption{Cleaning cost over annual recovered dollars, all 1{,}865 sites. No site reaches
  the breakeven line. Size targeting improves the ratio roughly $8\times$ and then
  asymptotes, because recovered dollars and cleaning cost both scale with panel count.}
  \label{fig:shortfall}
\end{figure*}

\paragraph{Scope of the verdict.}
IWSR is defined as the insolation-weighted mean of daily soiling ratios \emph{assuming
perfect cleaning}, so a permanent wash-only layer is excluded from the labels by
construction. Our loss model trains on $(1-\text{IWSR})\times100$ and therefore predicts
\emph{recoverable} loss, not total output deficit. Everything in this subsection is silent
about the other channel (\S\ref{sec:openchannel}).

\paragraph{The constant that was carrying the product.}
\texttt{recovery\_frac} was assumed at 0.90 and models at \textbf{0.045}, a $20\times$
overstatement. The 0.045 is a simulation output, not a measurement: it depends on a
fitted re-soiling constant and on the SOMOSclean saturation whose fit we qualify in
\S\ref{sec:unsourced}. The mechanism: at the calibrated re-soiling constant $k=15$
equivalent-days an array returns to its annual mean in about two weeks, and Santa Cruz
gets 27 heavy-rain ($\geq\SI{10}{\milli\metre}$) resets a year for free. One cleaning buys
roughly one array-month of clean out of twelve. Two validations: the same trajectory
reproduces NREL's measured coastal-California annual loss (model 5.06\% against measured
4.70\%), and the result is robust to $k$---even at $k=60$, four times the calibrated
value, in Phoenix, recovery is only 0.125. \textbf{With 0.90 restored, 44.8\% to 95.4\% of
the AOI ``should clean.''}

\paragraph{A counter-correction that does not move the verdict.}
Pricing every site at the current net-billing tariff was itself an untested assumption.
Public interconnection records~\cite{dgstats} filtered to the five AOI ZIP codes
($n=7{,}536$ residential PV) show \textbf{90.1\% of the AOI on legacy net metering}
(52.5\% NEM 2.0, 37.6\% NEM 1.0). The population-blended value of a lost kWh rises from
\$0.1646 to \$0.4284, a $2.60\times$ understatement, and the AOI annual
recoverable loss total rises 156\%. \textbf{The verdict did not move}: still zero of
1{,}865, still maximum $P(\text{net}>0) = 0.0000$.

\paragraph{Where the gap sits.}
Table~\ref{tab:shortfall} reports the shortfall by size band. Targeting by size buys
roughly $8\times$ and then asymptotes near $5\times$, because recovered dollars and
cleaning cost both scale with panel count. Small bands look worse only because the minimum
service charge dominates a twelve-panel job. Re-pricing across the published California
market range, \$5 to \$20 per panel, \textbf{zero sites pay at any price on the board}.

\begin{table}[t]
\centering\small
\caption{Shortfall (cleaning cost $\div$ annual recovered dollars) by system size. Values
$>1$ mean the wash costs more than it returns. Computed from
\artifact{site\_economics.csv} with the \emph{modelled} recovery fraction of 0.045 and the
level calibration applied. \S\ref{sec:results-econ} shows that fraction is denominator-
dependent (0.031 to 0.217) and reports the value of a wash in dollars instead; the
shortfall ordering below is unaffected.}
\label{tab:shortfall}
\begin{tabular}{lrrr}
\toprule
Size band (kW) & $n$ sites & Median & Best \\
\midrule
0--5      & 877 & $45.9\times$ & $26.5\times$ \\
5--10     & 796 & $24.6\times$ & $13.8\times$ \\
10--15    & 119 & $16.6\times$ & $13.7\times$ \\
15--30    & 35  & $13.8\times$ & $11.2\times$ \\
30--60    & 16  & $10.4\times$ & $8.6\times$ \\
60--150   & 12  & $8.5\times$  & $6.7\times$ \\
150+      & 10  & $5.5\times$  & $\mathbf{3.6\times}$ \\
\midrule
All       & \textbf{1{,}865} & --- & \textbf{0 clear} \\
\bottomrule
\end{tabular}
\end{table}

\subsection{Pilot deployment}
\label{sec:pilot}

The pipeline ran end to end on the AOI, producing 3{,}362 detected arrays resolving to
1{,}865 sites, a public dashboard with per-array size, orientation, and modelled loss, and
50 postcards rendered and mailed live on 2026-06-30 (50 of 50 accepted, \$47.52, QR
deep-links verified for all 50 identifiers). We present the mailing as a deliverability
and mechanics test, not as a behavioural experiment: a 50-card drop cannot resolve a
response rate statistically. The dashboard's honest claim is a per-home diagnosis: this is
your array, this is its size and orientation, this is what it loses, and no, do not pay to
clean it. \begin{figure}[t]
  \centering
  \includegraphics[width=0.62\columnwidth]{fig_deployment.pdf}
  \caption{The negative result as a homeowner receives it. The panel reports the loss in
  dollars, states that a wash leaves this roof about \$87 down, and says so first. It also
  discloses the channel we cannot measure---``a film only a wash removes\ldots{} nobody has
  measured how fast that builds up here''---and that we do not sell cleanings.
  \textbf{This capture predates the level calibration of \S\ref{sec:results-econ}}: the
  \SI{6.5}{\percent} shown is above anything the deployed data now contains, whose
  site-level losses have median \SI{2.77}{\percent} and maximum \SI{3.35}{\percent}. The
  overstatement runs in the product's favour and the panel still says do not pay, which is
  the property we are evidencing; the arithmetic on a current capture is smaller. The map
  is cropped out so no individual residence is identifiable.}
  \label{fig:deployment}
\end{figure}

Figure~\ref{fig:deployment} is that diagnosis, captured before the level calibration
described above. We show it because the claim
``the honest deliverable is telling someone not to buy'' is otherwise unverifiable, and
because the panel carries the paper's own caveat into the product: it tells the homeowner
that the figure assumes a best-week clean, that an average date recovers about a third as
much, and that a second soiling channel exists which nobody has measured here.

In the same footprint the previous-generation stack, built on
\SI{0.6}{\metre} NAIP imagery~\cite{naip}, found 334 arrays. Polygon counts are not
comparable across GSD (1.80 polygons per site at \SI{21}{\centi\metre} against 1.30 at
\SI{60}{\centi\metre}), sites are the comparable unit, and imagery and detector changed
together.

% ---------------------------------------------------------------------------
\section{Discussion}

\begin{table*}[t]
  \centering
  \small
  \caption{Every headline claim, the number behind it, and the strongest objection we know
  of. Limits are stated here rather than distributed through the text so that a reader can
  audit the paper's confidence in one place.}
  \label{tab:claims}
  \begin{tabular}{p{4.6cm}p{3.3cm}p{7.4cm}}
    \toprule
    \textbf{Claim} & \textbf{Number} & \textbf{Strongest limit we know of} \\
    \midrule
    Detection is gated and permissive
      & $F_1$ 0.826 [0.798, 0.853]
      & One county, one imagery vintage. Permit recall of 73.6\% is era-restricted; across
        all years it is 55.3\%. \\
    \addlinespace
    Our published temporal gate leaked
      & 88.7\% of held-out rows kept their own site in training
      & Measured, not inferred. This one has no serious objection; it is a property of the
        fold definition. \\
    \addlinespace
    Station labels cannot rank roofs
      & 0.622 [0.571, 0.670] jointly out-of-site and out-of-year, against 0.710 published
      & Robust to seed (sd 0.005, 7 seeds), to estimator (logistic regression gives 0.571),
        and to summary-row weighting (0.594--0.622). The panel label set is \SI{91}{\percent}
        western at station level; the east is censored rows only. \\
    \addlinespace
    The drop is leakage, not lost data
      & 83\% of the 0.0922 fall, 95\% CI [52, 99]\%
      & Controlled directly: a size-matched training set that keeps the leak scores 0.6986
        against 0.7144 full and 0.6222 clean. The joint fold does cost a median 28.8\% of
        the training rows, and we report that. \\
    \addlinespace
    Coordinates match the full feature set
      & 0.644 against 0.622
      & Not a significant win for coordinates either. The claim is the absence of evidence
        that 38 further features contribute, not evidence of their absence. \\
    \addlinespace
    Published physics models do no better
      & Kimber 0.610 [0.566, 0.652], unfitted, against our 0.622
      & Kimber's proxy is missing for 50 of 891 rows. The comparison is on ranking only;
        neither model was built to classify at our threshold. \\
    \addlinespace
    Does not travel to an unseen region
      & pooled 0.655, largest region 0.527
      & Stable across fold geometry (0.638--0.664 for $k=5\ldots12$). Regions are k-means
        cells, not a designed geographic sample. \\
    \addlinespace
    Per-system labels are necessary, not yet sufficient
      & levels 2.72 / 9.80 / 20.93 pts across three cleaning assumptions; rankings
        $\rho$ 0.241, 0.318, 0.894
      & Our own 0.573 uses \texttt{perfect\_clean}, which is the \emph{minority} ranker,
        and the only one whose level is plausible. The plausibility check is partly
        circular: NREL's IWSR is itself a perfect-clean construction. \\
    \addlinespace
    The dollar verdict is label-invariant
      & \$27.14 under all three assumptions
      & Algebraic, not empirical: annual loss cancels between the dollar figure and the
        recovery fraction's denominator. Holds only because we measure the step directly
        rather than inferring it from a modelled annual loss. \\
    \addlinespace
    One clean recovers \$28.10 of electricity
      & median over 505 events on 149 systems; \$150 service
      & Dry-window events are \emph{candidate} washes; an inverter reset or a shading change
        looks the same. Assumes cleaned and un-cleaned trajectories stay parallel. The
        \emph{fraction} of annual loss this represents is 0.031--0.217 depending on the
        cleaning assumption, which is why we report dollars. \\
    \addlinespace
    Cleaning does not clear
      & 0 of 149 at \$0.165/kWh, 2 of 149 at retail; median needs \$2.44/kWh; one wash
        returns \$28.10 at retail, \$10.14 at export
      & Recoverable soiling only; a permanent wash-only layer is excluded by construction.
        All 149 systems are Californian. Survives the \$150 service floor because the prize
        is smaller than any plausible price. \\
    \bottomrule
  \end{tabular}
\end{table*}

\label{sec:discussion}

\subsection{What the pipeline establishes}
Every stage runs on public data, and the answer the last stage returns is \emph{no}. That
is the result, not a defect. We reach \citet{mejia2013}'s conclusion for the average California site independently and
by a different route, but our claim is stronger and narrower. On recoverable soiling, at
sourced tariffs and a directly measured value per clean, cleaning never pays---not on the
1{,}865 detected sites in the pilot, and not on the 149 metered rooftops where the array
is instrumented and nothing about the loss has to be inferred from imagery. The second
population matters more than the first, because it removes detection, area estimation and
the risk model from the chain entirely and still returns the same answer.

\subsection{The open channel we cannot see}
\label{sec:openchannel}
Rain removes dust. Rain does not remove moss, lichen, algae, leaf litter, sap, or built-up
droppings, and shade plus moisture encourage them. That channel sits outside the IWSR
labels by construction, so nothing in \S\ref{sec:results-econ} speaks to it. We tested the
\emph{uniform} persistent-layer hypothesis on a revenue-grade PVDAQ site with 8.08 years
of data beside a class-A pyranometer: degradation $-0.076$\%/yr (95\% CI $[-0.408,
+0.358]$), standing layer 0.0 points, and the annual maximum soiling ratio pinned at
exactly 1.0000 in eight of eight years. That kills the uniform hypothesis. It does not
constrain the localised coastal case: that site is a ground mount beside farmland with no
canopy, and a whole-array ratio averages away a few bad modules.

The localised loss is wildly nonlinear and \emph{saturating}, not compounding: a shaded
cell current-limits its series string until a bypass diode cuts the substring out, so a
few square centimetres cost far more than their area fraction and then stop. Modelled at
cell granularity~\cite{bishop1988,holmgren2018pvlib}, four modules each losing one
substring cost 7.00\%, and a continuous edge band across four portrait modules on
string wiring costs 21.00\%, against breakeven bars of 2.2\% to 10.1\%. \textbf{A wash pays
here only if the roof carries continuous edge occlusion across several modules}. That is a
specific, falsifiable prediction about what a qualifying roof looks like, and it is
exactly what aerial imagery can see and a dust model cannot. Run the moss economics on
real per-roof geometry, however, and zero of 2{,}494
geometry-scored arrays clear a wash, best case losing about \$30, because module-level
electronics have been effectively mandated since 2019 and defuse the substring
nonlinearity. What remains
open is therefore the \emph{measurement}, not the product: nobody has measured biological
soiling's power cost in a Mediterranean climate. \citet{porcar2018} confirmed biofilm in
Berkeley but measured no power loss, and our AOI's growth-weighted time of wetness is
2{,}444\,h/yr against Berkeley's 2{,}468, a ratio of 0.99.

\subsection{A power result that reorders the obvious experiment}
A conventional before/after wash test on a single array saturates at roughly 5\% of output
minimum detectable effect and \emph{does not improve with more days}, because weather
arrives in multi-day blocks and autocorrelation rises at almost exactly the rate at which
extra samples would have helped. That is a floor measured on a site with an on-site pyranometer
and \SI{893}{\kilo\watt} of spatial averaging; a residential roof normalised against
satellite irradiance is $1.5$ to $2\times$ worse. The obvious study, wash 20 random homes and measure the average, is therefore both
underpowered and aimed at the wrong estimand. Selecting on visible soiling is not a shortcut; it is what makes the study
possible. A split-array wash on one instrumented volunteer roof reaches roughly 1.5\%
minimum detectable effect in two weeks for about \$150, because weather is common-mode and
cancels between halves.

\subsection{Limitations}
\label{sec:limitations}
\textbf{(a) Label transfer.} We train on instrumented ground stations and deploy on
residential rooftops. This is unvalidated and central.
\textbf{(b) Ranking.} \S\ref{sec:ranking}: valid for level, not for ranking homes.
\textbf{(c) Level.} The head reads high against every nearby measurement, and the scalar
that corrects it treats the symptom.
\textbf{(d) Recall.} 73.6\% on the imaged era and 55.3\% across all permit years; the
permit denominator is not a census.
\textbf{(e) Imagery staleness} is distinct from resolution. Any array installed after the
flight is invisible regardless of GSD. Our permit audit script still hard-codes a 2022
vintage constant against 2025 imagery, making the imaged-era split \emph{conservative};
this is a known defect, flagged, not yet fixed.
\textbf{(f) Unsourced constants}, listed in \S\ref{sec:unsourced}.
\textbf{(g) Inter-annotator agreement} was never measured, so the mask ceiling is unknown.
\textbf{(h) Geography.} Every metered rooftop in the economics is Californian, and the
panel labels behind the risk model are \SI{91}{\percent} western at station level
(Figure~\ref{fig:coverage}). The break-even boundary we map---$k\geq\SI{30}{\day}$, at
most 15 resets a year, above \$0.20/kWh---is therefore a \emph{prediction} about drier and
more expensive markets, not a measurement of them. It is the limitation we would close
first.
\textbf{(i) Label instability.} Per-system soiling labels are not yet well posed: three
standard cleaning assumptions place the median annual loss at 2.72, 9.80 and
\SI{20.93}{\percent} and rank the same systems at $\rho$ between 0.241 and 0.894, with the
only plausibly-levelled variant in the minority on ordering. This does not touch the dollar
result, which the cancellation of \S\ref{sec:results-econ} makes assumption-free, but it
does touch every ranking claim. Any programme that scales per-system label extraction should measure this before
trusting a ranking built on it.
\textbf{(j) Scale.} The free weather archive is the binding constraint on extending beyond
one region. In this work it cost roughly four hours of stalled fetching before we noticed
that most of what we were requesting was already cached under a different span.

\subsection{Unsourced constants, disclosed}
\label{sec:unsourced}
The minimum professional service charge (\$150) and the
per-panel schedule are our price points, not market quotes, and that one constant decides
most residential cases. Published California pricing runs \$5--12 per panel, so our
schedule sits at or below the bottom of the range and currently \emph{flatters} the
product. The production anchor of 1{,}500\,kWh/kWp/yr
is PVWatts-typical and unmeasured for this AOI. The packing factor (0.90) and midday
self-consumption share ($\sigma = 0.30$) are unsourced. The avoided-cost export table is
another utility's standing in for the local one, same methodology, level unverified. The SOMOSclean saturation parameter was fitted against the wrong side of the trajectory:
the objective scored 31 December, mid wet season, against an annual-mean measured IWSR.
That does not affect the annual-mean validation, but a refit would not be expected to
return 0.08. The net-metering cliff dates are encoded from memory and flagged in code as
requiring verification against the regulatory decision.

\subsection{Ethics and data provenance}
All inputs are public records used within their terms. Addresses are joined from county
parcel data to arrays detected in public imagery for the purpose of contacting the
property owner; we retain no personally identifying information beyond what the mailing
requires, and the public dashboard renders array geometry and modelled estimates only. We
state on the dashboard that these are modelled estimates from public data, not a
professional energy audit and not a guarantee. The deployed panel
(Figure~\ref{fig:deployment}) does report a percentage and a dollar figure for an
individual roof, which is a stronger claim than we would make in text, and we mark the
percentage as a county-level estimate rather than a measurement of that roof precisely
because we have no production data for it. We regard that labelling as the minimum, not as
sufficient: a per-address number carries an authority the underlying model does not earn,
and \S\ref{sec:ranking} is the reason why.

% ---------------------------------------------------------------------------
\section{Conclusion}
\label{sec:conclusion}

A soiling-targeting system can be built end to end on public data, and most of it works.
Arrays detect at $F_1$ 0.826 with 73.6\% independently measured recall, and per-array roof
geometry recovers from public lidar, agreeing with installer filings on the median tilt to
\SI{0.03}{\degree} while disagreeing about the distribution for a reason we can name.

What it cannot do is rank individual roofs, and the validation that suggested otherwise
was leaking: folds that hold out a site or a year but never both flatter the model by about
nine points. Corrected, the model scores 0.622, two coordinates score 0.644, and an
unfitted empirical formula from 2006 scores 0.610. Those three numbers say the same thing
from three directions. The cause is structural: public soiling labels are measured at
stations, so every roof in a catchment inherits one label and no per-roof feature has
anything to fit.

We tested the obvious repair rather than leaving it as an open question, and the answer is
partial. Per-system labels derived from public production telemetry \emph{do} carry
per-roof signal---0.573 within-cell concordance against a structural ceiling of 0.5---but
three standard ways of computing such a label disagree by a factor of eight on the level
and by as much as $\rho = 0.241$ on the ordering, and the one whose level is plausible is
the one the other two disagree with. The prerequisite for per-roof soiling targeting is therefore a label
estimator that is both unbiased for recoverable loss and tight enough to order two roofs in
one neighbourhood. That is a measurement problem, and it is the one worth working on.

Whether any of it would matter is settled separately, and more cheaply. On 149 metered
rooftops one wash recovers \$28.10 of electricity on the median system at the retail
offset, \$10.14 at the export rate, against a \$150 service; the median roof needs
electricity at \$2.44/kWh to break even. A wash must cost under \$10 before half of these
roofs clear at the export rate. No ranking model, however good, creates
value that is not there---which is why we would advise anyone entering this problem to
compute the arithmetic before building the pipeline, as we did not.

Two further channels remain genuinely open. The wash-only biological layer our labels
exclude by construction may carry measurable power loss in a Mediterranean climate; nobody
has measured it, and an afternoon with one instrumented roof would begin to settle it. And
every economic number here comes from Californian systems, so the boundary we map---$k \geq
\SI{30}{\day}$, at most 15 rain resets a year, above \$0.20/kWh---is a prediction about
drier and more expensive markets rather than a measurement of them.

% ---------------------------------------------------------------------------
\section*{Reproducibility}
\label{sec:repro}
Every number in this paper traces to a committed artifact, and the claim is checkable
rather than asserted: \texttt{paper/verify\_numbers.py} reads each headline value out of
the JSON that produced it and fails if the rounded form is absent from the manuscript. It
is wired into \texttt{make paper}, so the PDF cannot be built from stale claims. At
submission it passes 30 of 30, and it also fails the build if a \emph{retired} number
reappears.

The corrected validation results are produced by a standalone audit that imports nothing
from the pipeline it audits and reads the training matrix read-only, so a reader can
reproduce the negative results without trusting our code:

\begin{center}\small
\begin{tabular}{@{}ll@{}}
\toprule
Result & Script \\
\midrule
Fold taxonomy, leak measurement, 0.622 & \texttt{audit\_soiling\_validation.py} \\
Reporting-basis reconciliation         & \texttt{audit\_reconcile\_basis.py} \\
Out-of-region, year axis restored      & \texttt{audit\_regional\_holdout.py} \\
Seed, estimator, weighting, geometry   & \texttt{audit\_robustness.py} \\
Kimber and SOMOSclean baselines        & \texttt{audit\_physics\_baselines.py} \\
Cleaning events from telemetry         & \texttt{audit\_cleaning\_intervals.py} \\
Rain reset against manual wash         & \texttt{audit\_rain\_vs\_wash.py} \\
Measured recovery fraction             & \texttt{audit\_recovery\_empirical.py} \\
Within-cell ranking                    & \texttt{audit\_per\_system\_ranking.py} \\
Economics on metered rooftops          & \texttt{audit\_real\_systems\_econ.py} \\
Value per clean, label invariance      & \texttt{audit\_value\_per\_clean.py} \\
Three-assumption label comparison      & \texttt{audit\_label\_variants.py} \\
Training-size control                  & \texttt{audit\_trainsize\_control.py} \\
Break-even surface                     & \texttt{breakeven\_surface.py} \\
\midrule
Detection gate                         & \texttt{detect/eval\_tile\_f1.py} \\
Recall probe                           & \texttt{permit\_recall\_audit.py} \\
\bottomrule
\end{tabular}
\end{center}

\noindent
All of the above run from cached inputs and need no network. Every figure is emitted by a
script under \texttt{paper/figures/} that reads those artifacts directly; no figure
contains a hand-entered number.

\paragraph{What we got wrong, and how it was caught.} Three results in this paper changed
after they had been written down as findings, each because one setting was varied: the
recovery fraction turned out to have no single value at all---0.031, 0.063 or 0.217
according to which cleaning assumption supplies its denominator, which is why this paper
reports dollars per wash instead---the within-cell ranking fell from 0.573 to 0.512 on that
same choice, and the out-of-region figure moved once the production hyperparameters were
actually loaded. None
was found by review. We record this because it is the strongest argument we can make for
the robustness table in \S\ref{sec:results-risk} and for
\texttt{audit\_robustness.py} existing at all.

\section*{Acknowledgements}
We thank the Santa Cruz County GIS and building-permit offices, and NREL for publishing
both the soiling map and PVDAQ openly.

\bibliographystyle{plainnat}
\bibliography{refs}

\end{document}
